% !TEX root = ../matan.tex

\section{Непрерывные отображения, координатные функции.}

Рассмотрим отображение
\[
f\colon\Omega\subset\R^d\to\R^n,\qquad
f(x)=\bigl(f_1(x),\dots,f_n(x)\bigr),
\]
где $f_k\colon\Omega\to\R$ --- его \emph{координатные (компонентные) функции}, $f_k=\pi_k\circ f$, а $\pi_k\colon\R^n\to\R$ --- проекция на $k$-ю координату.

\begin{Definition}{Непрерывность отображения}{}
	Отображение $f\colon\Omega\subset\R^d\to\R^n$ \emph{непрерывно} в точке $x\in\Omega$, если
	\[
	\lim_{\substack{y\to x\\ y\in\Omega}}f(y)=f(x),
	\]
	то есть для любого $\eps>0$ существует $\delta>0$ такое, что из $y\in\Omega$, $|y-x|<\delta$ следует $|f(y)-f(x)|<\eps$ (нормы в $\R^d$ и $\R^n$ --- евклидовы).
\end{Definition}

\begin{Statement}{Покомпонентная непрерывность}{}
	Пусть $f=(f_1,\dots,f_n)\colon\Omega\subset\R^d\to\R^n$ и $x\in\Omega$. Тогда
	\[
	f\text{ непрерывна в }x
	\quad\Longleftrightarrow\quad
	f_k\text{ непрерывна в }x\ \text{ для всех }k=1,\dots,n.
	\]
\end{Statement}

\begin{proof}
	В основе --- двусторонняя оценка нормы вектора через его координаты: для любого $a=(a_1,\dots,a_n)\in\R^n$ и любого индекса $k$
	\[
	|a_k|\le\sqrt{\sum_{j=1}^{n}a_j^2}=|a|\le\sqrt{n}\,\max_{1\le j\le n}|a_j|.
	\]
	Применяем её к $a=f(y)-f(x)$, тогда $a_k=f_k(y)-f_k(x)$.

	($\Rightarrow$) Из левого неравенства $|f_k(y)-f_k(x)|\le|f(y)-f(x)|$: если правая часть стремится к нулю при $y\to x$, то и левая стремится к нулю, поэтому каждая $f_k$ непрерывна в $x$.

	($\Leftarrow$) Пусть все $f_k$ непрерывны в $x$. Зафиксируем $\eps>0$; для каждого $k$ найдётся $\delta_k>0$ такое, что $|f_k(y)-f_k(x)|<\eps$ при $y\in\Omega$, $|y-x|<\delta_k$. Положим $\delta=\min_{1\le k\le n}\delta_k>0$. Тогда при $|y-x|<\delta$ все координатные неравенства выполнены одновременно, и
	\[
	|f(y)-f(x)|=\sqrt{\sum_{k=1}^{n}|f_k(y)-f_k(x)|^2}\le\sqrt{n\eps^2}=\sqrt{n}\,\eps.
	\]
	Поскольку $\sqrt{n}\,\eps$ можно сделать сколь угодно малым, $f$ непрерывна в $x$.
\end{proof}

\begin{Remark}{Зачем нужны координатные функции}{}
	Утверждение сводит изучение непрерывности векторного отображения к изучению $n$ скалярных функций: $f$ непрерывна тогда и только тогда, когда непрерывна каждая её координата. Аналогичная редукция справедлива и для дифференцируемости (см. вопрос~38).
\end{Remark}

\begin{Theorem}{Непрерывность композиции}{}
	Пусть $f\colon\Omega\subset\R^d\to\R^n$ и $g\colon\widetilde\Omega\subset\R^n\to\R^m$, причём $f(\Omega)\subset\widetilde\Omega$. Если $f$ непрерывна в точке $x$, а $g$ непрерывна в точке $f(x)$, то композиция $g\circ f$ непрерывна в точке $x$.
\end{Theorem}

\begin{proof}
	Обозначим $u_0:=f(x)$ и зафиксируем $\eps>0$. По непрерывности $g$ в точке $u_0$ найдётся $\eta>0$ такое, что
	\[
	|g(u)-g(u_0)|<\eps\quad\text{при }u\in\widetilde\Omega,\ |u-u_0|<\eta.
	\]
	По непрерывности $f$ в точке $x$ для этого $\eta$ найдётся $\delta>0$ такое, что $|f(y)-f(x)|<\eta$ при $y\in\Omega$, $|y-x|<\delta$. Тогда при $|y-x|<\delta$ точка $u=f(y)$ лежит в $\widetilde\Omega$ и удовлетворяет $|u-u_0|<\eta$, откуда
	\[
	|g(f(y))-g(f(x))|<\eps.
	\]
	Значит, $g\circ f$ непрерывна в точке $x$.
\end{proof}

\begin{Theorem}{Равномерная непрерывность на компакте (теорема Кантора)}{}
	Если $f\colon K\to\R^n$ непрерывна на компакте $K\subset\R^d$, то $f$ равномерно непрерывна на $K$: для любого $\eps>0$ существует $\delta>0$ такое, что для всех $x,y\in K$ из $|x-y|<\delta$ следует $|f(x)-f(y)|<\eps$.
\end{Theorem}

\begin{proof}
	% [восстановлено]
	Предположим противное. Тогда найдётся $\eps_0>0$ и две последовательности $x_p,y_p\in K$ такие, что
	\[
	|x_p-y_p|<\frac1p,\qquad\text{но}\qquad |f(x_p)-f(y_p)|\ge\eps_0\quad\text{для всех }p.
	\]
	Множество $K$ компактно, поэтому из $\{x_p\}$ можно выделить сходящуюся подпоследовательность $x_{p_l}\to z\in K$. Так как $|x_{p_l}-y_{p_l}|\to0$, то и $y_{p_l}\to z$. По непрерывности $f$ в точке $z$ обе последовательности образов сходятся к $f(z)$:
	\[
	f(x_{p_l})\to f(z),\qquad f(y_{p_l})\to f(z),
	\]
	откуда $|f(x_{p_l})-f(y_{p_l})|\to0$. Это противоречит неравенству $|f(x_{p_l})-f(y_{p_l})|\ge\eps_0$. Значит, $f$ равномерно непрерывна на $K$.
\end{proof}
